% =====================================================================
% Kapitel 0: Vorwort
% =====================================================================

\chapter{Vorwort}

\section{Wofür dieses Handbuch}

Dieses Handbuch dokumentiert das Projekt \code{deephold\_db} -- eine
Finanzmarktdatenbank für Forschung und privaten Gebrauch. Es ist nicht
nur eine technische Referenz, sondern als \emph{Lehrbuch} geschrieben:
Es erklärt Konzepte, rechtfertigt Designentscheidungen und führt den
Leser Schritt für Schritt durch die Architektur.

Das Handbuch richtet sich an drei Zielgruppen:

\begin{itemize}
  \item \textbf{Entwickler}, die das Projekt verstehen, erweitern oder
        debuggen wollen. Sie brauchen einen vollständigen Überblick über
        Architektur, Datenmodell, ETL-Pipeline, Tests und Deployment.
  \item \textbf{Forscher}, die quantitative Finanzanalysen auf Basis
        der gesammelten Daten durchführen wollen -- etwa Walk-Forward-Backtests
        wie in Kapitel~\ref{ch:analytics} beschrieben.
  \item \textbf{Studierende und Einsteiger}, die ein Gefühl dafür
        bekommen wollen, wie ein real-existierendes Datenprodukt
        entsteht -- vom Docker-Container bis zur Performance-Metrik.
\end{itemize}

\begin{merke}
Dieses Handbuch setzt Programmierkenntnisse voraus -- idealerweise
Python und Grundkonzepte relationaler Datenbanken -- aber \emph{keine}
Finanzmarktexpertise. Alle Fachbegriffe werden bei ihrer ersten
Verwendung in einer \textbf{Definitionsbox} (blau) erklärt.
\end{merke}

\section{Was \code{deephold\_db} nicht ist}

\begin{warnung}
Dieses Projekt ist ausdrücklich \textbf{kein} Bloomberg- oder
Refinitiv-Ersatz. Es deckt bewusst nur Retail- und Research-Bedürfnisse ab:

\begin{itemize}
  \item \textbf{Keine} echten Point-in-Time-Daten -- die adjusted-close-Reihen
        von Yahoo Finance sind \emph{mit aktuellem Wissensstand korrigiert}
        (look-ahead-Bias bei echten Backtests vor 2010 möglich).
  \item \textbf{Keine} CRSP-äquivalente Aktienverlauf-Datenbank -- kein
        Delisting-Tracking, keine Survivorship-Bias-Korrektur.
  \item \textbf{Keine} Intraday- oder Tick-Daten -- alles ist
        Tagesfrequenz.
  \item \textbf{Keine} echte Continuous-Future-Konstruktion -- Front-Month
        Futures werden einzeln betrachtet, kein Roll-Kalender.
  \item \textbf{Keine} Option-Chains, keine Bond-CUSIPs, keine
        OTC-Derivate.
\end{itemize}
\end{warnung}

Für institutionelle Anwendungen mit Echtzeitdaten, vollständiger
PIT-Historie und Compliance-Audit-Trail ist dieses Projekt nicht
geeignet. Es ist explizit als \emph{Selbstlern- und Forschungsplattform}
konzipiert.

\section{Wie dieses Handbuch zu lesen ist}

\begin{beispiel}
\textbf{Lese-Empfehlungen}
  \begin{itemize}
    \item \textbf{Entwickler} lesen das Handbuch linear von vorne nach
          hinten. Die Kapitel bauen aufeinander auf.
    \item \textbf{Forscher} springen direkt zu Kapitel~\ref{ch:analytics}
          (AEGIS-Replikation) und nutzen Glossar (Anhang~A) zum
          Nachschlagen.
    \item \textbf{Einsteiger} beginnen mit Kapitel~\ref{ch:grundlagen}
          (Grundbegriffe) und lesen dann selektiv die Kapitel, die sie
          interessieren. Das Glossar hilft beim Verständnis der
          Fachbegriffe.
  \end{itemize}
\end{beispiel}

\section{Aufbau}

Das Handbuch gliedert sich in zehn Hauptkapitel und vier Anhänge:

\begin{description}
  \item[Kapitel 1 -- Einführung] Was das Projekt tut, seine Ziele, sein
        Scope und seine Anwendungsfälle.
  \item[Kapitel 2 -- Grundlagen] Erklärung aller Fachbegriffe, die im
        restlichen Handbuch vorkommen -- von PostgreSQL bis zu
        CAGR.
  \item[Kapitel 3 -- Infrastruktur] Docker, Docker Compose, das Makefile,
        Python-Environment und die Bun-Runtime.
  \item[Kapitel 4 -- Datenmodell] Das ER-Diagramm, die 14~Tabellen und
        ihre Indizes.
  \item[Kapitel 5 -- Vendor-Adapter] Die drei Datenlieferanten FRED, ECB
        und Yahoo Finance.
  \item[Kapitel 6 -- ETL-Pipeline] Seeder und das \code{query\_db.py}-Skript.
  \item[Kapitel 7 -- Notebook-Generator] Das Plotly-Notebook und sein
        Build-Prozess.
  \item[Kapitel 8 -- OpenTUI-Explorer] Das Terminal-UI auf Basis von Bun
        und React 19.
  \item[Kapitel 9 -- Tests] pytest, \code{bun test} und die manuelle
        Test-Checklist.
  \item[Kapitel 10 -- Analytics \& AEGIS-Replikation] Die
        Walk-Forward-Backtest-Implementierung und ihre Ergebnisse
        im Vergleich zum Originalpaper.
  \item[Anhang A -- Glossar] Alphabetisches Verzeichnis aller
        Fachbegriffe.
  \item[Anhang B -- Befehlsreferenz] Alle Makefile-Targets und
        CLI-Befehle.
  \item[Anhang C -- Datenbankschema] Vollständiges DDL aller Tabellen.
  \item[Anhang D -- Literatur] Bibliographie der zitierten Werke.
\end{description}

\section{Danksagung}

Dieses Projekt wäre ohne die folgenden Open-Source-Bibliotheken und
Datenquellen nicht möglich:

\begin{itemize}
  \item \textbf{PostgreSQL 16} -- die Datenbank-Engine.
  \item \textbf{SQLAlchemy 2.x} -- das ORM.
  \item \textbf{Pandas, Polars, NumPy} -- Datenverarbeitung in Python.
  \item \textbf{httpx, tenacity, structlog} -- robustes HTTP und Logging.
  \item \textbf{Prefect 2.x} -- Workflow-Orchestrierung (vorbereitet).
  \item \textbf{Plotly} -- interaktive Plots im Browser.
  \item \textbf{Bun, React 19, OpenTUI} -- das TUI-Frontend.
  \item \textbf{FRED, ECB, Yahoo Finance} -- die drei Datenlieferanten.
  \item \textbf{AEGIS-Paper} -- der Anlass für die Backtest-Replikation.
\end{itemize}

Vielen Dank an alle Maintainer dieser Projekte.
